%=============================================================================
% 模块五：调试、重构与安全 (约 8 页)
%=============================================================================

\section{调试与重构}

% -----------------------------------------------------------------------------
% 1. AI 辅助调试三部曲
% -----------------------------------------------------------------------------

\begin{frame}{AI 辅助调试 (Debug) 三部曲}
	\begin{center}
		\begin{tikzpicture}
			\node[draw, fill=red!20] (step1) at (0,0) {\shortstack{\textbf{第1步}\\提供 Traceback}};
			\node[draw, fill=yellow!20] (step2) at (4,0) {\shortstack{\textbf{第2步}\\说明环境}};
			\node[draw, fill=green!20] (step3) at (8,0) {\shortstack{\textbf{第3步}\\贴出数据}};
		\end{tikzpicture}
	\end{center}

	\vspace{0.3cm}

	\begin{columns}
		\column{0.33\textwidth}
		\textbf{第1步：提供 Traceback}

		完整复制报错信息：
		\begin{itemize}
			\item 错误类型
			\item 错误位置（行号）
			\item 完整的堆栈跟踪
		\end{itemize}

		\textbf{示例：}

		\texttt{cv2.error: ...}

		\texttt{Invalid number of channels}

		\column{0.33\textwidth}
		\textbf{第2步：说明环境}

		提供运行环境信息：
		\begin{itemize}
			\item 操作系统
			\item Python 版本
			\item OpenCV 版本
			\item 其他相关库版本
		\end{itemize}

		\textbf{示例：}

		\texttt{Windows 11}

		\texttt{Python 3.9.7}

		\texttt{OpenCV 4.8.0}

		\column{0.33\textwidth}
		\textbf{第3步：贴出输入数据}

		提供输入数据信息：
		\begin{itemize}
			\item 数据类型和格式
			\item 数据的维度/大小
			\item 期望的输出格式
			\item 实际的输出（如果有错误）
		\end{itemize}

		\textbf{示例：}

		\texttt{图像尺寸：1920x1080}

		\texttt{通道：3 (RGB)}
	\end{columns}
\end{frame}

\begin{frame}{调试 Prompt 模板}
	\begin{exampleblock}{完整的调试 Prompt 示例}
		我的 OpenCV 代码报错了，请帮我分析：

		\vspace{0.2cm}

		\textbf{错误信息：}

		\texttt{cv2.error: OpenCV(4.8.0) ...}

		\texttt{Invalid number of channels in input image}

		\vspace{0.2cm}

		\textbf{运行环境：}

		\texttt{Windows 11, Python 3.9.7, OpenCV 4.8.0}

		\vspace{0.2cm}

		\textbf{相关代码：}

		\texttt{img = cv2.imread('test.jpg')}

		\texttt{gray = cv2.cvtColor(img, cv2.COLOR\_BGR2GRAY)}

		\vspace{0.2cm}

		\textbf{输入数据：}

		\texttt{图像文件 test.jpg，尺寸 1920x1080}

		\vspace{0.2cm}

		\textbf{期望行为：} 正常转换为灰度图
	\end{exampleblock}
\end{frame}

% -----------------------------------------------------------------------------
% 2. 重构：把 AI 当成导师
% -----------------------------------------------------------------------------

\begin{frame}{重构：把 AI 当成导师}
	\begin{block}{重构的意义}
		代码能运行只是第一步，\textbf{优雅的代码}才能长期维护。
	\end{block}

	\vspace{0.3cm}

	\begin{columns}
		\column{0.5\textwidth}
		\textbf{重构 Prompt 模板：}

		\begin{exampleblock}{重构请求}
			请帮我重构这段代码，提高代码质量：

			要求：
			\begin{enumerate}
				\item 提高运行效率
				\item 添加类型提示
				\item 改进命名规范
			\end{enumerate}
		\end{exampleblock}

		\column{0.5\textwidth}
		\textbf{AI 可能的重构建议：}

		\begin{enumerate}
			\item \textbf{性能优化}：向量化运算、避免重复计算
			\item \textbf{可读性}：变量名有意义、函数拆分合理
			\item \textbf{健壮性}：输入校验、异常处理
		\end{enumerate}
	\end{columns}
\end{frame}

% -----------------------------------------------------------------------------
% 3. AI工具使用规范：鼓励的操作
% -----------------------------------------------------------------------------

\begin{frame}{AI工具使用规范：鼓励的操作}
	\begin{exampleblock}{我们强烈鼓励以下行为}
		\begin{enumerate}
			\item \textbf{设计高质量的Prompt}
			\begin{itemize}
				\item 使用RTF框架（Role-Task-Format）
				\item 记录成功的Prompt模板
				\item 分享有效的Prompt给同学
			\end{itemize}

			\item \textbf{用AI学习新的库和API}
			\begin{itemize}
				\item 请AI解释官方文档
				\item 请求代码示例和最佳实践
				\item 理解库的设计思想
			\end{itemize}
		\end{enumerate}
	\end{exampleblock}

	\vspace{0.3cm}

	\begin{block}{学习提示}
		AI是你的\textbf{学习伙伴}，而不是\textbf{代写工具}。善用AI可以提高学习效率，但不能替代你的思考过程。
	\end{block}
\end{frame}

\begin{frame}{更多鼓励的行为}
	\begin{exampleblock}{继续鼓励的行为}
		\begin{enumerate}
			\item \textbf{用AI分析代码逻辑}
			\begin{itemize}
				\item 代码审查，发现潜在bug
				\item 分析算法复杂度
				\item 理解他人代码
			\end{itemize}

			\item \textbf{记录学习过程}
			\begin{itemize}
				\item 保存与AI的有价值对话
				\item 整理常见问题及解决方案
				\item 建立个人知识库
			\end{itemize}
		\end{enumerate}
	\end{exampleblock}

	\vspace{0.3cm}

	\begin{block}{为什么要记录？}
		\begin{itemize}
			\item 好的Prompt可以被复用
			\item 解决问题的思路可以积累
			\item 形成个人的知识体系
		\end{itemize}
	\end{block}
\end{frame}

% -----------------------------------------------------------------------------
% 4. 安全与伦理
% -----------------------------------------------------------------------------

\begin{frame}{安全与伦理：红线警告}
	\begin{alertblock}{\faExclamationTriangle\ 重要提醒}
		使用 AI 辅助编程时，必须注意以下安全与伦理问题！
	\end{alertblock}

	\vspace{0.3cm}

	\begin{columns}
		\column{0.5\textwidth}
		\textbf{1. 数据安全}

		\begin{block}{不要上传}
			\begin{itemize}
				\item API Key、密码
				\item 个人隐私数据
				\item 商业机密代码
			\end{itemize}
		\end{block}

		\begin{block}{安全做法}
			\begin{itemize}
				\item 使用脱敏示例数据
				\item 本地部署模型
			\end{itemize}
		\end{block}

		\column{0.5\textwidth}
		\textbf{2. 版权与责任}

		\begin{block}{代码归属}
			AI生成代码的版权复杂，检查所用工具的许可协议
		\end{block}

		\begin{block}{责任界定}
			AI生成的代码可能有错误，最终责任人是你自己
		\end{block}
	\end{columns}
\end{frame}

\begin{frame}{学术诚信}
	\begin{block}{学术诚信要求}
		\begin{itemize}
			\item 作业/论文中使用 AI 需声明
			\item 了解学校的 AI 使用政策
			\item 保持学习和思考的过程
		\end{itemize}
	\end{block}

	\vspace{0.3cm}

	\begin{alertblock}{核心原则}
		\begin{itemize}
			\item AI 是\textbf{辅助工具}，不是\textbf{替代品}
			\item 学习过程比结果更重要
			\item 理解代码原理是期末考试的重点
		\end{itemize}
	\end{alertblock}
\end{frame}
